% !TEX root = ../main.tex
\clearpage
\appendix
\chapter{평가 파이프라인 및 재현성}
\label{ch:appendix-eval}

본 부록은 \ref{ch:results}장 수치를 재현하는 데 사용된 공개 평가 파이프라인을 문서화한다. 저장소 구조, 오프라인 합성 평가, 실제 BigQuery 추출 절차, 보조 정렬 요약 표를 포함한다.

\dissertationSubheading[sec:repo-layout]{저장소 구조}

평가 파이프라인은 \texttt{A-Skill-Programs/margin\_rank/}에 있다. (1) 최신 ERC-20 allowance, (2) 타임스탬프 송금, (3) 디코딩된 GMX \texttt{PositionDecrease} 이벤트를 수집한다. EndorseRank, AWP, GF-PR 점수는 \texttt{compute\_reputation\_ranks.py}를 통해 \texttt{wallet\_rankings.parquet}에 병합된다. 다섯 프록시 계열은 \texttt{proxy\_metrics.py}에서 계산; 정렬 및 벤치마크는 \texttt{run\_dissertation\_eval.py}를 통해 실행.

주요 출력: \texttt{data/processed/eval\_summary.json}(수치) 및 \texttt{2-Dissertation-Draft/results/tables/} 하 LaTeX 조각. 확장 baseline(7-method, 6-Aave)은 \texttt{2-Dissertation-Draft/archive/extended-baselines/}로 내보낸다.

\dissertationSubheading[sec:offline-eval]{오프라인 합성 평가 (Google Cloud 불필요)}

\begin{verbatim}
cd A-Skill-Programs/margin_rank
python3 -m venv .venv && source .venv/bin/activate
pip install -r requirements.txt
python scripts/run_dissertation_eval.py --fixtures --n-wallets 571 --export-latex
cd ../../2-Dissertation-Draft && ./build.ps1
\end{verbatim}

\dissertationSubheading[sec:bq-extraction]{실제 BigQuery 추출}

Application Default Credentials 및 \texttt{GOOGLE\_CLOUD\_PROJECT} 필요. tier 및 비용 가드레일은 \texttt{docs/bigquery\_data\_plan.md} 참조.

\begin{verbatim}
export GOOGLE_CLOUD_PROJECT=your-project-id

# Phase 1 — trading-success label extraction (GMX V2) (6-month window)
python scripts/extract_margin_week.py --dry-run
python scripts/extract_margin_week.py --extract --yes
python scripts/decode_gmx_events.py
python scripts/compute_rankings.py

# Phase 2 — Reputation subgraph (monthly batches)
python scripts/extract_reputation_data.py --dry-run
python scripts/extract_reputation_data.py --extract --yes
python scripts/preprocess_arbitrum_allowances.py
python scripts/preprocess_arbitrum_transfers.py
python scripts/compute_reputation_ranks.py

# Phase 3 — Five-proxy evaluation + Tier 2/3 (local; no extra BQ for robustness)
python scripts/sample_active_wallets.py --from-parquet --extract
python scripts/run_dissertation_eval.py --real --robustness --export-latex

# Optional Tier-2 BQ re-extract (preserves Tier-1 raw under data/raw/reputation/)
python scripts/extract_reputation_data.py --tier2 --dry-run
python scripts/extract_reputation_data.py --tier2 --extract --yes --resume
python scripts/preprocess_tier2_reputation.py
\end{verbatim}

추출 메타데이터(처리 바이트, 행 수)는 \texttt{data/processed/extraction\_manifest.json}에 기록.

\dissertationSubheading[sec:tier2-tier3-appendix]{Tier~2 및 Tier~3 표}

표~\ref{tab:benchmark-scaling-incohort}는 Tier~1 in-cohort 런타임 스케일링($n \leq 5{,}521$)을 유지한다. Tier~3 강건성 표는 \ref{ch:results}장에 있으며 \texttt{run\_dissertation\_eval.py --real --robustness --export-latex}로 재생성한다.

% Real BigQuery data — regenerate with:
%   A-Skill-Programs/margin_rank/scripts/run_dissertation_eval.py --real --export-latex
% Generated by export_latex_results.py — do not edit by hand unless re-export fails

\begin{table}[htbp]
\centering
\caption{In-cohort runtime scaling (SHA256 seed \texttt{benchmark-scale-v1}; $n \leq 5{,}521$ matched GMX cohort).}
\label{tab:benchmark-scaling-incohort}
\begin{tabular}{rrrrrr}
\toprule
$n$ & ER runtime (s) & AWP runtime (s) & ER speedup (AWP/ER) & ER edges & AWP edges \\
\midrule
1000 & 0.393 & 3.215 & 8.2$\times$ & 2584 & 23058 \\
2000 & 0.729 & 6.794 & 9.3$\times$ & 5147 & 47031 \\
4000 & 2.075 & 18.031 & 8.7$\times$ & 10557 & 97105 \\
5521 & 2.105 & 18.711 & 8.9$\times$ & 14727 & 137087 \\
\bottomrule
\end{tabular}
\end{table}


\dissertationSubheading[sec:delegation-deferred]{Credit Delegation 추출 (연기)}

Aave V3 \texttt{BorrowAllowanceDelegated} 로그는 \texttt{margin\_config.yaml}에 구성되어 있으나 본 연구에서 추출하지 않았다(\texttt{aave\_delegation\_events.parquet} 부재). 보관된 six-Aave baseline(\texttt{Delegation-PR})은 따라서 degenerate 정렬($\tau = \text{---}$)을 보고한다. 후속 추출 패스에서 전용 delegation SQL 템플릿과 \texttt{extract\_aave\_lending.py --include-delegation}으로 재활성화.

\dissertationSubheading[sec:supplementary-summary]{보조 정렬 요약 표}

표~\ref{tab:alignment-summary}는 표~\ref{tab:alignment-three-methods}의 계열 수준 평균 $\tau$를 열 형식으로 중복(EndorseRank / AWP / GF-PR / 최고 정렬 레이블). CSV 지향 재현성을 위해 유지; 본문은 표~\ref{tab:alignment-three-methods}만 사용.

% DUMMY DATA — Replace after running:
%   A-Skill-Programs/margin_rank/scripts/run_dissertation_eval.py --fixtures --n-wallets 571
% Real data: BigQuery goog_blockchain_arbitrum_one_us + GMX EventEmitter logs
% Generated by export_latex_results.py — do not edit by hand unless re-export fails

\begin{table}[htbp]
\centering
\caption{Cross-proxy mean Kendall $\tau$: transfer family vs.\ GMX margin family (Claim~3--4 summary).}
\label{tab:alignment-summary}
\begin{tabular}{lcc}
\toprule
Method & Transfer proxies (mean $\tau$) & GMX proxies (mean $\tau$) \\
\midrule
EndorseRank & 0.466 & 0.506 \\
AWP & 0.289 & 0.466 \\
\bottomrule
\end{tabular}
\end{table}

